% ===== BOT 17e: dao sau muc 17.5 -- no N khong tran, scale_power, sqrt(12) =====

\begin{frame}{Đào sâu 17.5 (1/3): rủi ro nổ $N$ khi không có trần cho $k$}
\footnotesize
Mục 17.5 đã nêu: \texttt{gaussian\_model.py:698} (comment) ghi ``clamp min=2 và max=8'' nhưng
\texttt{:701} (code thật) chỉ \texttt{ks = torch.clamp(ks, min=1)} --- \textbf{không có trần trên}.
Ở đây định lượng mức độ nổ khi giả định cả 3 trục cùng vi phạm bằng $\eta$ (worst case đối xứng):
\[
k=\Big\lceil\sqrt{\max(\eta,1)}\Big\rceil,\qquad N=k_xk_yk_z
\]
\vspace{0.1cm}
\begin{center}
\includegraphics[width=0.92\linewidth]{deep17e_01_N_explosion_no_clamp.png}
\end{center}
\begin{itemize}
  \item $\eta{=}(9,9,9)\Rightarrow N{=}27$ (ví dụ đã nêu ở 17.5); tới $\eta{=}64$ code thật và bản
        ``có trần max=8'' còn trùng nhau --- lệch nhau chỉ lộ rõ từ $\eta{>}64$: tại $\eta{=}100$,
        code thật cho $N{=}1000$ so với $N{=}512$ nếu có trần --- gần \textbf{gấp đôi}.
  \item $\eta$ lớn nhất \textbf{đo được thật} trong checkpoint (G2, trục $z$) $=491.3\Rightarrow k_z{=}23$
        trên \emph{một} trục; $N$ thật của G2 chỉ là $184$ (2 trục còn lại $\eta$ nhỏ) --- bằng chứng số
        cho tính dị hướng: $N$ chỉ nổ theo trục thật sự vi phạm.
\end{itemize}
\end{frame}


\begin{frame}{Đào sâu 17.5 (2/3): \texttt{scale\_power} $p$ đổi kích thước, không đổi vị trí}
\footnotesize
Trên Gaussian $G_1$ \textbf{thật} của checkpoint \texttt{point\_cloud\_iter5000.ply}
($s_{\text{old}}{=}0.53156$, $k{=}3$ thật trên trục $z$), cố định $k$, so sánh $p\in\{1.0,1.5,2.0\}$:
\[
\boxed{\,s_{\text{new}}=\dfrac{s_{\text{old}}}{k^{\,p}}\,}\quad(\texttt{gaussian\_model.py:722})
\qquad\qquad
\text{offset}=\dfrac{s_{\text{old}}}{k}\sqrt{12}\cdot\text{grid}\ \ (\texttt{:775--778, không phụ thuộc }p)
\]
\vspace{0.1cm}
\begin{center}
\includegraphics[width=0.92\linewidth]{deep17e_02_scale_power_effect.png}
\end{center}
\begin{itemize}
  \item Khoảng cách giữa 2 con liền kề \textbf{cố định} $=\sqrt{12}\,(s_{\text{old}}/k)=0.6138$ (không đổi theo $p$).
  \item $p{=}1.0$: khe hở $=0.2594$; $p{=}1.5$: khe hở $=0.4092$; $p{=}2.0$: khe hở $=0.4957$ ---
        \textbf{gần gấp đôi} từ $p{=}1$ lên $p{=}2$, vì offset không co lại cùng tốc độ với kích thước con.
\end{itemize}
\end{frame}


\begin{frame}{Đào sâu 17.5 (3/3): vì sao $\sqrt{12}$, không phải $\sqrt3$ hay $\sqrt6$}
\footnotesize
Suy luận đúng theo code: \texttt{separations} là bề rộng \textbf{một ô lưới} (nửa bề rộng $a=\text{sep}/2$),
$\mathrm{Var}(\mathcal U(-a,a))=a^2/3$ phải khớp $\sigma_{\text{new}}^2$:
\[
\boxed{\,\dfrac{(\text{sep}/2)^2}{3}=\sigma_{\text{new}}^2 \;\Longleftrightarrow\; \text{sep}=\sigma_{\text{new}}\sqrt{12}\,}
\qquad(\texttt{gaussian\_model.py:775--776})
\]
\vspace{0.1cm}
\begin{center}
\includegraphics[width=0.78\linewidth]{deep17e_03_sqrt12_grid_variance.png}
\end{center}
\begin{itemize}
  \item Trên $G_1$ thật ($k{=}3$, $\sigma_{\text{new}}^2{=}0.031395$): chỉ $\sqrt{12}$ cho
        $\mathrm{Var}(\text{1 ô lưới})=0.031395$ khớp \textbf{chính xác}; $\sqrt3\Rightarrow0.007849$
        (nhỏ hơn 4 lần), $\sqrt6\Rightarrow0.015698$ (nhỏ hơn 2 lần).
  \item \alert{Sắc thái chưa nêu ở 17.5:} phương sai \textbf{thực nghiệm} của $k{=}3$ điểm lưới rời rạc
        lại lớn hơn mục tiêu theo hệ số $(k^2{-}1){=}8$ ở cả 3 lựa chọn --- $\sqrt{12}$ chỉ khớp đúng
        phương sai giả định cho \textbf{một ô lưới}, không phải đẳng thức chính xác cho toàn bộ $k$ điểm.
\end{itemize}
\end{frame}
